% Part: second-order-logic
% Chapter: metatheory
% Section: introduction

\documentclass[../../../include/open-logic-section]{subfiles}

\begin{document}

\olfileid[hi]{sol}{met}{int}

\olsection{परिचय}

प्रथम-कोटि तर्कशास्त्र में कई उपयोगी गुण हैं। हम जानते हैं कि यह
निर्णेय नहीं है, किंतु कम-से-कम यह अभिगृहीतीकरणीय है। अर्थात् प्रथम-कोटि
तर्कशास्त्र के लिए ऐसी प्रमाण-पद्धतियाँ हैं जो ध्वनि और पूर्ण हैं; दूसरे
शब्दों में, वे ऐसा !!a{derivability} संबंध~$\Proves$ देती हैं जिसका गुण यह
है कि !!{sentence} के किसी भी समुच्चय~$\Gamma$ और !!{sentence}~$!A$ के लिए
$\Gamma \Entails !A$ तभी और केवल तभी जब $\Gamma \Proves !A$। विशेषतः
इसका अर्थ है कि प्रथम-कोटि तर्कशास्त्र की वैधताएँ !!{computably
  enumerable} हैं। कोई संगणनीय फलन~$f\colon \Nat \to \Sent[L]$ ऐसा है
कि~$f$ के मान ठीक वही और केवल वही !!{sentence} हैं जो~$\Lang{L}$ में वैध
हैं। इसका कारण यह है कि !!{derivation} का परिगणन किया जा सकता है, और जो
!!{derive} किसी एक~!!{sentence} को करते हैं उन्हें फिर उसी !!{sentence} पर
प्रतिचित्रित किया जा सकता है। द्वितीय-कोटि तर्कशास्त्र प्रथम-कोटि
तर्कशास्त्र से अधिक अभिव्यंजक है, इसलिए सामान्यतः उसकी वैधताओं को निरूपित
करना अधिक जटिल है। वास्तव में, हम दिखाएँगे कि द्वितीय-कोटि तर्कशास्त्र न
केवल अनिर्णेय है, बल्कि उसकी वैधताएँ !!{computably enumerable} भी नहीं हैं।
इसका अर्थ है कि द्वितीय-कोटि तर्कशास्त्र के लिए कोई ध्वनि और पूर्ण
प्रमाण-पद्धति नहीं हो सकती (यद्यपि ध्वनि पर अपूर्ण प्रमाण-पद्धतियाँ उपलब्ध
हैं और वास्तव में शोध की महत्त्वपूर्ण विषय-वस्तुएँ हैं)।

प्रथम-कोटि तर्कशास्त्र में दो और गुण भी हैं: वह संहत है (यदि !!{sentence}
के किसी समुच्चय~$\Gamma$ का प्रत्येक परिमित उपसमुच्चय संतुष्ट्य है, तो
$\Gamma$ स्वयं संतुष्ट्य है), और उसके लिए ल्योवेनहाइम--स्कोलेम प्रमेय लागू
होता है (यदि $\Gamma$ का कोई अनंत मॉडल है, तो उसका कोई !!a{denumerable}
मॉडल भी है)। ये दोनों परिणाम द्वितीय-कोटि तर्कशास्त्र के लिए विफल होते हैं।
यहाँ भी कारण यही है कि द्वितीय-कोटि तर्कशास्त्र !!{domain} के आकार के विषय
में वे तथ्य व्यक्त कर सकता है जिन्हें प्रथम-कोटि तर्कशास्त्र व्यक्त नहीं
कर सकता।

\end{document}
